\chapter{Making Decisions: Control Flow and Pattern Matching}

\section{The Fork in the Road: Conditionals}

\begin{kidbox}{Choosing What to Wear}
Every morning, you look out the bedroom window.
\begin{itemize}
    \item \textbf{If} it is raining outside, you put on your yellow raincoat and boots!
    \item \textbf{Else if} it is freezing cold and snowing, you put on your puffy winter coat and mittens!
    \item \textbf{Else} (it is sunny!), you put on your favorite dinosaur t-shirt and sunglasses!
\end{itemize}
In Shae, computers make choices using \texttt{if}, \texttt{else if}, and \texttt{else}:
\begin{lstlisting}[language=shaelang]
let weather = "rainy";

if weather == "rainy" {
    println("Grab your yellow raincoat!");
} else if weather == "snowy" {
    println("Put on your warm mittens!");
} else {
    println("Time for dinosaur sunglasses!");
}
\end{lstlisting}
\end{kidbox}

\begin{syntaxbox}{If Expressions in Shae}
In Shae, \texttt{if} statements are true expressions: they return a value directly!
\begin{lstlisting}[language=shaelang]
let age = 10;
let status = if age >= 18 { "Adult" } else { "Child" };
println("Status: " + status); // "Child"
\end{lstlisting}
Both branches must evaluate to compatible expressions when assigned to a variable.
\end{syntaxbox}

\section{Loops: The Merry-Go-Round}

\begin{kidbox}{Spinning on the Carousel}
Have you ever ridden a carousel at the fair?
You get on a painted wooden horse. The music plays, and the carousel spins around and around!
You keep spinning \textbf{while} the music is playing!
When the music stops, you step off the horse.
And if you want to say hello to every friend in line, you can visit them one by one!
\begin{lstlisting}[language=shaelang]
// 1. While Loop (Music is playing!)
let spins = 0;
while spins < 3 {
    println("Wheee! Spin number " + spins);
    spins = spins + 1;
}

// 2. For .. in Loop (Visiting friends in line)
let friends = ["Alice", "Bob", "Charlie"];
for friend in friends {
    println("Hello, " + friend + "!");
}
\end{lstlisting}
\end{kidbox}

\begin{deepbox}{Branch Prediction and Loop Compilation}
At the assembly and bytecode level, loops are compiled into a conditional forward branch (\texttt{OpJumpIfFalse}) and an unconditional backwards jump (\texttt{OpLoop}):
\begin{enumerate}
    \item Condition evaluated onto stack.
    \item \texttt{OpJumpIfFalse exit\_label}: Checks condition. If falsy, jumps past loop body.
    \item Loop body statements executed.
    \item \texttt{OpLoop start\_label}: Computes negative relative jump offset and rewinds the instruction pointer \texttt{ip}.
\end{enumerate}
CPUs optimize loops via the Branch Target Buffer (BTB). Because backwards jumps almost always evaluate as taken (until the final iteration), hardware branch predictors achieve $>98\%$ branch prediction accuracy.
\end{deepbox}

\section{Pattern Matching: The Ultimate Shape Sorter}

\begin{kidbox}{The Magical Toy Sorter}
Imagine a magical sorting machine with special slots.
You drop a mystery box in. The machine checks:
"Is it a number between 1 and 10? Is it a secret word? Is it a dinosaur with wings?"
In Shae, \texttt{match} is the ultimate decision maker:
\begin{lstlisting}[language=shaelang]
let dice_roll = 6;

match dice_roll {
    1 => println("You rolled a 1! Try again!"),
    2..=5 => println("Nice middle roll!"),
    6 => println("Critical hit! Maximum score!"),
    _ => println("Invalid dice!"),
}
\end{lstlisting}
\end{kidbox}

\begin{syntaxbox}{Match Guards and Range Patterns}
Shae pattern matching supports value matching, wildcard fallbacks, range patterns, and boolean match guards:
\begin{lstlisting}[language=shaelang]
let temperature = 28;

match temperature {
    temp if temp < 0 => println("Freezing ice!"),
    0..15 => println("Chilly autumn breeze"),
    15..=25 => println("Perfect room temperature"),
    temp if temp > 35 => println("Heatwave warning!"),
    _ => println("Warm and pleasant"),
}
\end{lstlisting}
\end{syntaxbox}
